% ΣΗΜΜΥ — source exported in HTML-only mode. Compile with XeLaTeX.
\documentclass[a4paper,11pt]{article}
\usepackage{fontspec,polyglossia,amsmath,amssymb,geometry,xcolor,graphicx}
\setmainlanguage{greek}
\IfFontExistsTF{DejaVu Serif}{\setmainfont{DejaVu Serif}}{\setmainfont{FreeSerif}}
\geometry{margin=23mm}
\definecolor{accent}{HTML}{7446A7}
\setlength{\parindent}{0pt}
\begin{document}
\begin{titlepage}\centering\vspace*{2cm}
{\Large\color{accent}\bfseries ΟΜΑΔΑ ΕΠΙΛΥΣΗΣ ΘΕΜΑΤΩΝ ΣΗΜΜΥ\par}\vspace{2cm}
{\Huge\bfseries Γραμμική Άλγεβρα\par}\vspace{0.7cm}
{\Large Κανονική Εξέταση -- 2025\par}\vfill
\begin{tabular}{rl}Λύτης: &  \\
Ελεγκτής Α: &  \\
Ελεγκτής Β: & \end{tabular}
\vfill
{\itshape Λύσεις γραμμένες από φοιτητές για φοιτητές\par}\end{titlepage}
\section*{Άσκηση 1 α)}

\textbf{Μεθοδολογία}\par

\begin{enumerate}
\item Υπολογίζουμε την ισοδύναμη αντίσταση.
\item Εφαρμόζουμε τον νόμο του Ohm.
\end{enumerate}

\textbf{Λύση}\par

\[
R_{\mathrm{eq}} = R_1+R_2 = 20\,\Omega
\]

\[
I=\frac{100}{20}=5\,\mathrm{A}
\]

Επομένως, το ρεύμα είναι \(I=5\,\mathrm{A}\).\par

\section*{Άσκηση 1 β)}

\textbf{Μεθοδολογία}\par

Χρησιμοποιούμε το αποτέλεσμα του α) ερωτήματος.\par

\textbf{Λύση}\par

Η ισχύς προκύπτει από την σχέση \(P=I^2R\).\par

\end{document}
